% GENERATED FILE. Edit canonical lessons, then run npm run generate:books.
% canonical-source-hash: fnv1a64:143cf37c3ce90704
% canonical-lessons: PA-C275-dhilla, PA-C275-tang, PA-C275-phabbna, PA-C275-darjan, PA-C275-thaila

\chapter{Too Loose, Too Tight}
\label{ch:pa-a2-too-loose-too-tight}
\hlchaptermodality{275}

\begin{chapteropening}
\textbf{By the end of this chapter:} \emph{I can say whether clothes fit me.}

Complete the last lesson of chapter 275: I can say whether clothes fit me.
\end{chapteropening}

\section[\texorpdfstring{\pa{ਢਿੱਲਾ} (ḍhillā)}{ਢਿੱਲਾ (ḍhillā)}]{\pa{ਢਿੱਲਾ} (ḍhillā) — loose; slack}

\label{lesson:PA-C275-dhilla}



\begin{quote}
Before the new one: say the Punjabi for a size, then the Punjabi for a pair.
\end{quote}

\subsection*{You'll want to know: \pa{ਢਿੱਲਾ}}
\textbf{\pa{ਢਿੱਲਾ}} — \emph{ḍhillā} — \textquotedblleft{}loose; slack\textquotedblright{}.

\begin{quote}
\pa{ਇਹ} \pa{ਕੁੜਤਾ} \pa{ਮੈਨੂੰ} \pa{ਢਿੱਲਾ} \pa{ਹੈ।}
\end{quote}

\begin{quote}
\emph{eh kuṛtā mainū̃ ḍhillā hai.}
\end{quote}

\begin{quote}
This kurta is loose on me.
\end{quote}

\begin{cousinweb}
\textbf{\pa{ਢਿੱਲਾ}} (\emph{ḍhillā}) is inherited from a reconstructed \emph{*ḍhila}, \textquotedblleft{}loose, slack\textquotedblright{}; CDIAL lists Punjabi \emph{ḍhillā}, \textquotedblleft{}slack, slow, late\textquotedblright{}, under it (CDIAL 5590).
\end{cousinweb}

\begin{grammarlens}[title={What you've built}]
One more word for this chapter.
\end{grammarlens}

\subsection*{Your turn}
\begin{itemize}
\raggedright
  \item \emph{Say it:} \emph{ḍhillā}
  \item \emph{Say it:} \emph{ḍhillā}, once more
  \item \emph{Say it:} \emph{eh kuṛtā mainū̃ ḍhillā hai.} — This kurta is loose on me.
  \item \emph{Recall:} say the Punjabi for a size, then the Punjabi for a pair, then say \emph{ḍhillā} again
\end{itemize}

\begin{tcolorbox}[breakable,colback=teal!4,colframe=teal!35!black,title={Before you move on}]
What does \pa{ਢਿੱਲਾ} mean? (\textquotedblleft{}Loose; slack\textquotedblright{}.) Say it once more.
\end{tcolorbox}
\section[\texorpdfstring{\pa{ਤੰਗ} (taṅg)}{ਤੰਗ (taṅg)}]{\pa{ਤੰਗ} (taṅg) — tight}

\label{lesson:PA-C275-tang}



\begin{quote}
Before the new one: say the Punjabi for a pair, then the Punjabi for loose.
\end{quote}

\subsection*{You'll want to know: \pa{ਤੰਗ}}
\textbf{\pa{ਤੰਗ}} — \emph{taṅg} — \textquotedblleft{}tight\textquotedblright{}.

\begin{quote}
\pa{ਇਹ} \pa{ਜੁੱਤੀ} \pa{ਬਹੁਤ} \pa{ਤੰਗ} \pa{ਹੈ।}
\end{quote}

\begin{quote}
\emph{eh juttī bahut taṅg hai.}
\end{quote}

\begin{quote}
This shoe is very tight.
\end{quote}

\begin{cousinweb}
\textbf{\pa{ਤੰਗ}} (\emph{taṅg}) is a loan, usually traced to Persian \emph{tang}, \textquotedblleft{}narrow, tight\textquotedblright{}.
\end{cousinweb}

\begin{grammarlens}[title={What you've built}]
One more word for this chapter.
\end{grammarlens}

\subsection*{Your turn}
\begin{itemize}
\raggedright
  \item \emph{Say it:} \emph{taṅg}
  \item \emph{Say it:} \emph{taṅg}, once more
  \item \emph{Say it:} \emph{eh juttī bahut taṅg hai.} — This shoe is very tight.
  \item \emph{Recall:} say the Punjabi for a pair, then the Punjabi for loose, then say \emph{taṅg} again
  \item \emph{Write it:} copy \textbf{\pa{ਤੰਗ}} once, checking the whole word
\end{itemize}

\begin{tcolorbox}[breakable,colback=teal!4,colframe=teal!35!black,title={Before you move on}]
What does \pa{ਤੰਗ} mean? (\textquotedblleft{}Tight\textquotedblright{}.) Say it once more.
\end{tcolorbox}
\section[\texorpdfstring{\pa{ਫੱਬਣਾ} (phabbṇā)}{ਫੱਬਣਾ (phabbṇā)}]{\pa{ਫੱਬਣਾ} (phabbṇā) — to suit; to look good on someone}

\label{lesson:PA-C275-phabbna}



\begin{quote}
Before the new one: say the Punjabi for loose, then the Punjabi for tight.
\end{quote}

\subsection*{You'll want to know: \pa{ਫੱਬਣਾ}}
\textbf{\pa{ਫੱਬਣਾ}} — \emph{phabbṇā} — \textquotedblleft{}to suit; to look good on someone\textquotedblright{}.

\begin{quote}
\pa{ਨੀਲਾ} \pa{ਕੁੜਤਾ} \pa{ਅਮਨ} \pa{ਨੂੰ} \pa{ਬਹੁਤ} \pa{ਫੱਬਦਾ} \pa{ਹੈ।}
\end{quote}

\begin{quote}
\emph{nīlā kuṛtā Aman nū̃ bahut phabbdā hai.}
\end{quote}

\begin{quote}
The blue kurta really suits Aman.
\end{quote}

Past (it …): \pa{ਉਹ} \pa{ਕੁੜਤਾ} \pa{ਅਮਨ} \pa{ਨੂੰ} \pa{ਫੱਬਿਆ}, \emph{oh kuṛtā Aman nū̃ phabbiā}.

\begin{cousinweb}
\textbf{\pa{ਫੱਬਣਾ}} (\emph{phabbṇā}) is inherited from a reconstructed \emph{*sparvati}, \textquotedblleft{}releases, succeeds, is fitting\textquotedblright{}; CDIAL lists Punjabi \emph{phabbṇā}, \textquotedblleft{}to look well, be profitable\textquotedblright{}, under it (CDIAL 13808).
\end{cousinweb}

\begin{grammarlens}[title={What you've built}]
One more word for this chapter.
\end{grammarlens}

\subsection*{Your turn}
\begin{itemize}
\raggedright
  \item \emph{Say it:} \emph{phabbṇā}
  \item \emph{Say it:} \emph{phabbṇā}, once more
  \item \emph{Say it:} \emph{nīlā kuṛtā Aman nū̃ bahut phabbdā hai.} — The blue kurta really suits Aman.
  \item \emph{Recall:} say the Punjabi for loose, then the Punjabi for tight, then say \emph{phabbṇā} again
  \item \emph{Write it:} copy \textbf{\pa{ਫੱਬਣਾ}} once, checking the whole word
\end{itemize}

\begin{tcolorbox}[breakable,colback=teal!4,colframe=teal!35!black,title={Before you move on}]
What does \pa{ਫੱਬਣਾ} mean? (\textquotedblleft{}To suit; to look good on someone\textquotedblright{}.) Say it once more.
\end{tcolorbox}
\section[\texorpdfstring{\pa{ਦਰਜਨ} (darjan)}{ਦਰਜਨ (darjan)}]{\pa{ਦਰਜਨ} (darjan) — a dozen}

\label{lesson:PA-C275-darjan}



\begin{quote}
Before the new one: say the Punjabi for tight, then the Punjabi for to suit.
\end{quote}

\subsection*{You'll want to know: \pa{ਦਰਜਨ}}
\textbf{\pa{ਦਰਜਨ}} — \emph{darjan} — \textquotedblleft{}a dozen\textquotedblright{}.

\begin{quote}
\pa{ਮੈਂ} \pa{ਇੱਕ} \pa{ਦਰਜਨ} \pa{ਕੇਲੇ} \pa{ਖ਼ਰੀਦਾਂਗਾ।}
\end{quote}

\begin{quote}
\emph{maiṁ ikk darjan kele kharīdāṅgā.}
\end{quote}

\begin{quote}
I will buy a dozen bananas.
\end{quote}

A masculine noun.

\begin{cousinweb}
\textbf{\pa{ਦਰਜਨ}} (\emph{darjan}) is borrowed from English \emph{dozen}.
\end{cousinweb}

\begin{grammarlens}[title={What you've built}]
One more word for this chapter.
\end{grammarlens}

\subsection*{Your turn}
\begin{itemize}
\raggedright
  \item \emph{Say it:} \emph{darjan}
  \item \emph{Say it:} \emph{darjan}, once more
  \item \emph{Say it:} \emph{maiṁ ikk darjan kele kharīdāṅgā.} — I will buy a dozen bananas.
  \item \emph{Recall:} say the Punjabi for tight, then the Punjabi for to suit, then say \emph{darjan} again
  \item \emph{Write it:} copy \textbf{\pa{ਦਰਜਨ}} once, checking the whole word
\end{itemize}

\begin{tcolorbox}[breakable,colback=teal!4,colframe=teal!35!black,title={Before you move on}]
What does \pa{ਦਰਜਨ} mean? (\textquotedblleft{}A dozen\textquotedblright{}.) Say it once more.
\end{tcolorbox}
\section[\texorpdfstring{\pa{ਥੈਲਾ} (thailā)}{ਥੈਲਾ (thailā)}]{\pa{ਥੈਲਾ} (thailā) — a bag; a shopping bag}

\label{lesson:PA-C275-thaila}



\begin{quote}
Before the new one: say the Punjabi for to suit, then the Punjabi for a dozen.
\end{quote}

\subsection*{You'll want to know: \pa{ਥੈਲਾ}}
\textbf{\pa{ਥੈਲਾ}} — \emph{thailā} — \textquotedblleft{}a bag; a shopping bag\textquotedblright{}.

\begin{quote}
\pa{ਮੇਰਾ} \pa{ਥੈਲਾ} \pa{ਸਬਜ਼ੀ} \pa{ਨਾਲ} \pa{ਭਰਿਆ} \pa{ਹੈ।}
\end{quote}

\begin{quote}
\emph{merā thailā sabzī nāl bhariā hai.}
\end{quote}

\begin{quote}
My bag is full of vegetables.
\end{quote}

A masculine noun.

\begin{cousinweb}
\textbf{\pa{ਥੈਲਾ}} (\emph{thailā}) is inherited from Sanskrit \emph{sthavi}, \textquotedblleft{}sack, bag\textquotedblright{}; CDIAL lists Punjabi \emph{thailā}, \textquotedblleft{}bag\textquotedblright{}, under it (CDIAL 13746).
\end{cousinweb}

\begin{grammarlens}[title={What you've built}]
One more word for this chapter.
\end{grammarlens}

\subsection*{Your turn}
\begin{itemize}
\raggedright
  \item \emph{Say it:} \emph{thailā}
  \item \emph{Say it:} \emph{thailā}, once more
  \item \emph{Say it:} \emph{merā thailā sabzī nāl bhariā hai.} — My bag is full of vegetables.
  \item \emph{Recall:} say the Punjabi for to suit, then the Punjabi for a dozen, then say \emph{thailā} again
  \item \emph{Write it:} copy \textbf{\pa{ਥੈਲਾ}} once, checking the whole word
\end{itemize}

\begin{tcolorbox}[breakable,colback=teal!4,colframe=teal!35!black,title={Before you move on}]
What does \pa{ਥੈਲਾ} mean? (\textquotedblleft{}A bag; a shopping bag\textquotedblright{}.) Say it once more.
\end{tcolorbox}
